\ExerciseSection

\begin{enumerate}
\def\labelenumi{\arabic{enumi})}
\item
  Let X be a Hausdorff uniform space and let \((f_{t})\) be a family of pseudometrics on X which define the uniformity of X. Let \(X_{t}\) be the metric space associated with the uniform space obtained by endowing X with the structure defined by the single pseudometric \(f_{t}\) (no. 1). Show that X is isomorphic to a subspace of the product uniform space \(\prod_{i} X_{i}\) .
\item
  In a connected metric space X for which the metric is not bounded on \(X \times X\) , show that a sphere cannot be empty.
\end{enumerate}

¶ 3) a) Let X be a compact metric space. If the closure in X of every open ball is the closed ball of the same centre and radius, then every ball in X is a connected set {[}if x and y are two points of X, and if S is the closed ball with centre x and radius \(d(x, y)\) , show that for each \(\varepsilon > 0\) the set \(A_{y,\varepsilon}\) of points of S which can be joined to y by a \(V_{\varepsilon}\) -chain contained in S (Chapter II, § 4, no. 4) contains points z such that \(d(x, z) < d(x, y)\) {]}. Deduce that \(x \in A_{y,\varepsilon}\) for all \(\varepsilon > 0\) , arguing by contradiction and using Weierstrass's theorem (Chapter IV, § 6, no. 1, Theorem 1).

\begin{enumerate}
\def\labelenumi{\alph{enumi})}
\setcounter{enumi}{1}
\tightlist
\item
  In the space \(\mathbf{Y} = \mathbf{R}^2\) with the metric
\end{enumerate}

\[
d (x, y) = \max \left(\left| x _ {1} - y _ {1} \right|, \left| x _ {2} - y _ {2} \right|\right),
\]

let X be the compact subspace consisting of all points \((x_{1}, x_{2})\) such that either \(x_{1} = 0\) and \(0 \leqslant x_{2} \leqslant 1\) , or \(0 \leqslant x_{1} \leqslant 1\) and \(x_{2} = 0\) . Show that every ball in X is connected but that the closure of an open ball is not necessarily the closed ball with the same centre and radius.

\begin{enumerate}
\def\labelenumi{\arabic{enumi})}
\setcounter{enumi}{3}
\tightlist
\item
  A metric space \(\mathbf{X}\) is said to be an ultrametric space if its metric \(d\) satisfies the inequality
\end{enumerate}

\[
d (x, y) \leqslant \sup \left(d (x, z), d (y, z)\right)
\]

for all \(x, y, z\) in \(\mathbf{X}\) (which implies the triangle inequality) (cf.~§ 6, Exercise 2).

\begin{enumerate}
\def\labelenumi{\alph{enumi})}
\item
  If \(d(x, z) \neq d(y, z)\) , show that \(d(x, y) = \sup (d(x, z), d(y, z))\) .
\item
  Let \(\mathrm{V}_{r}(x)\) be the open ball with centre x and radius r. Show that \(\mathrm{V}_{r}(x)\) is both open and closed in X (and hence that X is totally disconnected) and that for each \(y \in \mathrm{V}_{r}(x)\) we have \(\mathrm{V}_{r}(y) = \mathrm{V}_{r}(x)\) .
\item
  Show that the closed ball \(\mathrm{W}_{r}(x)\) with centre x and radius r is both open and closed in X, and that for each \(y \in \mathrm{W}_{r}(x)\) we have
\end{enumerate}

\[
\mathrm{W} _ {r} (y) = \mathrm{W} _ {r} (x).
\]

The distinct open balls of radius r contained in \(W_{r}(x)\) form a partition of \(W_{r}(x)\) , and the distance between any pair of these balls is equal to r. d) If two balls (open or closed) of X intersect, then one is contained in the other.

\begin{enumerate}
\def\labelenumi{\alph{enumi})}
\setcounter{enumi}{4}
\item
  A sequence \((x_{n})\) of points of X is a Cauchy sequence if and only if \(d(x_{n}, x_{n+1})\) tends to 0 as \(n\) tends to infinity.
\item
  If X is compact, show that for each \(x_{0} \in X\) the set of values of \(d(x_{0}, x)\) in X is a (finite or infinite) countable subset of \([0, +\infty]\) , of which all the points, with the possible exception of 0, are isolated {[}for each value r taken by \(d(x_{0}, x)\) , consider the least upper bound of \(d(x_{0}, x)\) on the set of points where \(d(x_{0}, x) < r\) , and its greatest lower bound on the set of points where \(d(x_{0}, x) > r\) {]}.
\end{enumerate}

\begin{enumerate}
\def\labelenumi{\arabic{enumi})}
\setcounter{enumi}{4}
\tightlist
\item
  Let X be a metric space, d its metric. For each pair \((x, y)\) of points of X, let \(d_{0}(x, y)\) denote the greatest lower bound of the real numbers \(\alpha > 0\) such that x and y can be joined by a \(V_{\alpha}\) -chain (Chapter II, § 4, no. 4). Show that \(d_{0}\) is a pseudometric on X, and that the metric space associated with the uniform space defined by the pseudometric \(d_{0}\) on X (no. 1) is an ultrametric space (Exercise 4).
\end{enumerate}

§ 6) Let X be a metric space. If A and B are two non-empty subsets of X, put

\[
\rho (\mathrm{A}, \mathrm{B}) = \sup _ {x \in \mathbf {A}} d (x, \mathrm{B}) \quad \text { and } \quad \sigma (\mathrm{A}, \mathrm{B}) = \sup \left(\rho (\mathrm{A}, \mathrm{B}), \rho (\mathrm{B}, \mathrm{A})\right);
\]

also put \(\sigma(\emptyset, \emptyset) = 0\) , \(\sigma(\emptyset, A) = \sigma(A, \emptyset) = +\infty\) for any non-empty subset \(A\) of \(X\) . Show that \(\sigma\) is a pseudometric on the set \(\mathfrak{P}(X)\) of all subsets of \(X\) , and that the uniformity defined by \(\sigma\) coincides with the uniformity constructed from that of \(X\) by the procedure of Chapter II, § 1, Exercise 5.

Suppose that X is bounded. Then \(\sigma\) is a metric on the set \(\mathfrak{F}(\mathbf{X})\) of non-empty closed subsets of X. If, moreover, X is complete, show that \(\mathfrak{F}(\mathbf{X})\) is a complete metric space {[}let \(\Phi\) be a Cauchy filter on \(\mathfrak{F}(\mathbf{X})\) ; for each set \(\mathcal{X} \in \Phi\) , let \(S(\mathcal{X})\) be the union of the subsets A of X which belong to \(\mathcal{X}\) ; show that the sets \(S(\mathcal{X})\) form a filter base on X as \(\mathcal{X}\) runs through \(\Phi\) , that this filter base has a non-empty set C of cluster points, and that \(\Phi\) converges to C{]}.

\begin{enumerate}
\def\labelenumi{\arabic{enumi})}
\setcounter{enumi}{6}
\item
  Let X be an infinite discrete space. Show that the uniformity of finite partitions (Chapter II, § 2, no. 2) on X, which is compatible with the topology of X, is not metrizable (otherwise there exists a sequence \((\mathfrak{F}_n)\) of finite partitions of X such that every finite partition of X is formed of sets, each of which is a union of sets belonging to one of the partitions \(\mathfrak{F}_n\) ; from this it would follow that the set of finite partitions of X is countable).
\item
  Let \((\mathbf{X}_{t})\) be an uncountable family of Hausdorff topological spaces, each of which has at least two distinct points. Show that, in the product space \(\prod_{i} X_{i}\) , no point has a countable fundamental system of neighbourhoods.
\item
  Let X be a topological space, each point of which has a countable fundamental system of neighbourhoods.
\end{enumerate}

\begin{enumerate}
\def\labelenumi{\alph{enumi})}
\item
  X is Hausdorff provided that every convergent sequence in X has only one limit.
\item
  If a is a cluster point of an infinite sequence \((x_{n})\) of points of X, then there is an infinite subsequence of \((x_{n})\) which converges to a.
\item
  Let A be a non-empty subset of X, let \(x_{0}\) be a point of \(\overline{A}\) , and let f be a mapping of A into a Hausdorff topological space \(X'\) . Then a point \(a \in X'\) is a limit of f at the point \(x_{0}\) , relative to A, provided that, for every sequence \((x_{n})\) of points of A which converges to \(x_{0}\) , the sequence \((f(x_{n}))\) converges to a.
\item
  With the notation of c), suppose also that every point of \(X'\) has a countable fundamental system of neighbourhoods. If \(a \in X'\) is a cluster point of f at \(x_{0}\) relative to A, then there is a sequence \((x_{n})\) of points of A which converges to \(x_{0}\) and is such that the sequence \((f(x_{n}))\) converges to a.
\end{enumerate}

¶ 10) Let X be a compact space. If there exists a continuous real-valued function f on X×X such that the relation \(f(x,y)=0\) is equivalent to x=y, then X is metrizable. (Show that if V runs through a fundamental system of neighbourhoods of 0 in R, the sets \(\overline{f}^{1}(V)\) form a fundamental system of neighbourhoods of the diagonal \(\Delta\) in X×X.)

¶ 11) Let X be a compact metric space, with metric d.~Show that if f is a mapping of X into X such that, for each pair x, y of points of X, we have \(d(f(x), f(y)) \geqslant d(x, y)\) , then f is an isometry of X onto itself. {[}Let a, b be any two points of X, and put \(f^n = f^{n-1} \circ f\) , \(a_n = f^n(a)\) , \(b_n = f^n(b)\) ; show that, for each \(\varepsilon > 0\) , there exists an index k such that \(d(a, a_k) \leqslant \varepsilon\) and \(d(b, b_k) \leqslant \varepsilon\) , by choosing suitable subsequences of \((a_n)\) and \((b_n)\) ; hence show that \(d(a_1, b_1) = d(a, b)\) and that \(f(\mathbf{X})\) is dense in X.{]}

¶ 12) Let X be a compact metrizable space.

\begin{enumerate}
\def\labelenumi{\alph{enumi})}
\item
  Show that there is a continuous mapping of Cantor's triadic set K (Chapter IV, § 2, no. 5) onto X (cf.~Chapter IV, § 8, Exercise 11).
\item
  If in addition X is totally disconnected and has no isolated points, then X is homeomorphic to K. {[}Argue as in Chapter IV, § 8, Exercise 12, using the fact that every neighbourhood of a point \(x \in \mathbf{X}\) contains a neighbourhood of \(x\) which is both open and closed (Chapter II, § 4, no. 4, Corollary to Proposition 6){]}.
\end{enumerate}

¶ 13) a) On the set R of real numbers, consider the topology C defined as follows: for each y \textgreater{} 0, \(\mathrm{U}_{y}(x)\) denotes the union of the intervals \([x, x + y[ \text{ and } ] - x - y, -x[; C \text{ is the topology for which the } \mathrm{U}_{y}(x) \text{ form a fundamental system of neighbourhoods of } x \text{ as } y \text{ runs through the set of real numbers } > 0. \text{ Let X be the space obtained by endowing the interval } [-1, +1] \text{ with the topology induced by C. Show that X is compact. (Consider a filter F on X; if } x \in X \text{ is a cluster point of F with respect to the topology of the real line, show that either } x \text{ or } -x \text{ is a cluster point of F with respect to C.)}\)

\begin{enumerate}
\def\labelenumi{\alph{enumi})}
\setcounter{enumi}{1}
\item
  Every point of X has a countable fundamental system of neighbourhoods, and X has a countable dense subset, but X is not metrizable (show that its topology has no countable base).
\item
  Let A be an open subset of X. Show that A is the union of a countable family \((I_{n})\) of open intervals contained in {[}0, 1{]}, the intervals \(-I_{n}\) , a subset of the set of left-hand points of the intervals \(I_{n}\) and \(-I_{n}\) , and possibly the point +1. Hence show that A is the union of a countable family of closed subsets of X.
\item
  Let Y be the set \(I \times \{1, 2\}\) , where I is the interval {[}0, 1{]} of R; let \(X_{i}\) denote \(I \times \{i\} (i = 1, 2)\) and let f denote the bijection \((x, 1) \to (x, 2)\) of \(X_{1}\) onto \(X_{2}\) . For each \(x \in I\) , let \(\mathfrak{B}((x, 2))\) denote the set of subsets consisting only of \(\{(x, 2)\}\) , and let \(\mathfrak{B}((x, 1))\) denote the set of subsets of the form \(V_{1} \cup (f(V_{1}) - \{x, 2\})\) , where \(V_{1} = V \times \{1\}\) and V runs through a fundamental system of neighbourhoods of x in I. Show that for each \(y \in Y\) , \(\mathfrak{B}(y)\) is a fundamental system of neighbourhoods of y for a topology on Y. Endowed with this topology, Y is compact and every point of Y has a countable fundamental system of neighbourhoods, but Y has no countable dense subset. Every closed subset of Y contained in \(X_{2}\) is finite; hence show that the compact subset \(X_{1}\) of Y has no countable fundamental system of neighbourhoods. If R is the equivalence relation on Y whose classes are the set \(X_{1}\) and the points of \(Y - X_{1}\) , R is closed and every equivalence class mod R is compact, but there is a point in Y/R which has no countable fundamental system of neighbourhoods {[}cf.~§4, Exercise 24a){]}.
\end{enumerate}

\begin{enumerate}
\def\labelenumi{\arabic{enumi})}
\setcounter{enumi}{13}
\tightlist
\item
  \begin{enumerate}
  \def\labelenumii{\alph{enumii})}
  \tightlist
  \item
    Let X be an accessible topological space (Chapter I, § 8, Exercise 1). Show that the following properties are equivalent:
  \end{enumerate}
\end{enumerate}

α) Every sequence of points of X has a cluster point.

\(\beta)\) No infinite discrete subspace of \(\mathbf{X}\) is closed.

\(\gamma)\) Every countable open covering of \(\mathbf{X}\) contains a finite open covering \(\mathbf{X}\) .

δ) Given any infinite open covering R of X, there exists an open covering S ⊂ R of X, distinct from R.

A topological space X is said to be countably compact if it is Hausdorff and has the above properties.

\begin{enumerate}
\def\labelenumi{\alph{enumi})}
\setcounter{enumi}{1}
\item
  A sequence of points in a countably compact space is convergent if and only if it has a unique cluster point.
\item
  Every closed subspace of a countably compact space is countably compact. Conversely, if X is Hausdorff and if every point of X has a countable fundamental system of neighbourhoods, then every countably compact subspace of X is closed in X.
\item
  Let \(f\) be a continuous mapping of a countably compact space \(\mathbf{X}\) into a Hausdorff space \(\mathbf{X}'\) . Then \(f(\mathbf{X})\) is a countably compact subspace of \(\mathbf{X}'\) .
\item
  Let X be a countably compact space, every point of which has a countable fundamental system of neighbourhoods. Then every sequence of points of X has a convergent subsequence.
\end{enumerate}

\(f\) ) Let \((\mathbf{X}_n)\) be a countable sequence of topological spaces, in each of which every point has a countable fundamental system of neighbourhoods.

The product space \(\prod_{n=1}^{\infty} \mathbf{X}_n\) is then countably compact if and only if each of the spaces \(\mathbf{X}_n\) is countably compact. {[}Use e) and Chapter I, §6, Exercise 16.{]}

\begin{enumerate}
\def\labelenumi{\alph{enumi})}
\setcounter{enumi}{6}
\item
  A countably compact space in which every point has a countable fundamental system of neighbourhoods is regular.
\item
  If a countably compact space has a countable base of open sets, it is compact and therefore metrizable.
\item
  Every lower semi-continuous real-valued function on a countably compact space attains its greatest lower bound. In particular, every countably compact space is pseudo-compact (§ 1, Exercise 21) {[}cf.~§ 4, Exercise 26 b){]}.
\item
  Show that the property \(\delta\) of \(a\) implies the following:
\end{enumerate}

ζ) Every point-finite (§ 4, no. 3) open covering of X contains a finite open covering of X. (Argue by contradiction.) Conversely, a regular space which satisfies ζ) is countably compact {[}show that ζ) then implies β){]}.

¶ 15) Let \(\mathbf{X} = [a, b[\) be the locally compact space defined in Chapter I, § 9, Exercise 12.

\begin{enumerate}
\def\labelenumi{\alph{enumi})}
\item
  A subset of X is relatively compact if and only if it is bounded. Hence show that X is countably compact and therefore (Proposition 15) non-metrizable (observe that every countable subset of X is bounded).
\item
  Show that every point of X has a metrizable neighbourhood (use Proposition 16).
\item
  If A and B are two non-compact closed sets in X, their intersection is not empty (form an increasing sequence in which the even-numbered terms are points of A and the odd-numbered terms are points of B). Deduce that every neighbourhood of a non-compact closed set is the complement of a relatively compact set. If A is the set of non-isolated points of X, show that A is closed and is not the intersection of any countable family of open sets.
\item
  Let \(\mathbf{X}'\) denote the interval \([a, b]\) endowed with the following topology: for each \(x \in \mathbf{X}\) the intervals \(]y, x]\) , where \(y\) runs through the set of elements \(< x\) , form a fundamental system of neighbourhoods of \(x\) ; a fundamental system of neighbourhoods of \(b\) is formed by the sets \(\mathbf{V}_x\) where, for each \(x \in \mathbf{X}\) , \(\mathbf{V}_x\) denotes the union of \(\{b\}\) and the set of points of \([x, b]\) which have a predecessor (i.e., which are isolated in \(\mathbf{X}\) ). Show that \(\mathbf{X}'\) is countably compact but not regular {[}cf.~Exercise 14 g){]} and that the subspace \(\mathbf{X}\) of \(\mathbf{X}'\) is countably compact but not closed {[}cf.~Exercise 14 c){]}.
\end{enumerate}

\begin{enumerate}
\def\labelenumi{\arabic{enumi})}
\setcounter{enumi}{15}
\tightlist
\item
  Let X and Y be two countably compact spaces. Show that the product \(X \times Y\) is countably compact in each of the following two cases: (i) one of X, Y is compact; (ii) one of X, Y is such that each point has a countable fundamental system of neighbourhoods. {[}In case (i), supposing that Y is compact, consider an increasing sequence \((\mathbf{G}_{n})\) of open sets whose union is \(X \times Y\) ; for each n, let \(H_{n}\) be the set of all \(x \in X\) such that \(\{x\} \times Y \subset G_{n}\) ; show that \(H_{n}\) is open and that \(X = \bigcup_{n} H_{n} \cdot\) {]}
\end{enumerate}

¶ 17) Let βN be the Stone-Cech compactification of the discrete space N (§ 1, Exercise 7).

\begin{enumerate}
\def\labelenumi{\alph{enumi})}
\item
  Show that there are two countably compact subspaces A, B of \(\beta N\) such that \(A \cap B = N\) and \(A \cup B = \beta N\) . {[}Let \(\aleph_{\alpha} = 2^{2^{\operatorname{Card}(N)}}\) ; by virtue of § 1, Exercise 12 b), there is a bijection \(\xi \to S_{\xi}\) of the ordinal \(\omega_{\alpha}\) (Set Theory, Chapter III, § 6, Exercise 10) onto the set of countably infinite subsets of X. By transfinite induction, define two injective mappings \(\xi \to x_{\xi}, \xi \to y_{\xi}\) of \(\omega_{\alpha}\) into \(\beta N - N\) such that \(x_{\xi} \in \overline{S}_{\xi}, y_{\xi} \in \overline{S}_{\xi}\) for each \(\xi\) , and such that if P (resp. Q) is the set of the \(x_{\xi}\) (resp. \(y_{\xi}\) ) we have \(P \cap Q = \varnothing\) , \(P \cup Q = \beta N - N\) ; for this, make use of § 1, Exercise 12 b) and c). Then show that A = P ∪ N and B = Q ∪ N satisfy the conditions of the question.{]}
\item
  Show that the product space \(A \times B\) is not countably compact (note that the intersection of \(A \times B\) with the diagonal of \(\beta N \times \beta N\) is an infinite discrete closed subspace of \(A \times B\) ).
\end{enumerate}

¶ 18) Let X be a metric space such that, for each \(x \in X\) , there is an open ball with centre \(x\) which, when considered as a subspace of X, has a countable base. Let \(r_x\) be the least upper bound of the radii of balls with centre \(x\) which have this property.

\begin{enumerate}
\def\labelenumi{\alph{enumi})}
\item
  Show, using Zorn's lemma, that there is a maximal family \((\mathbf{B}_{\alpha})\) of open balls in X which are pairwise disjoint and such that, if \(x_{\alpha}\) is the centre of \(B_{\alpha}\) , the radius of \(B_{\alpha}\) is \(<r_{x_{\alpha}}\) . Show that the union of the \(B_{\alpha}\) is dense in X, and hence that there is a dense subset M of X such that every open ball with centre at an arbitrary point \(x \in X\) and radius \(<r_{x}\) contains only a countable infinity of points of M (note that such a ball can meet only a countable infinity of pairwise disjoint open sets, and use Proposition 12).
\item
  For each point \(x \in \mathbf{M}\) , let \(\mathbf{S}_x\) denote the open ball with centre \(x\) and radius \(\frac{1}{3} r_x\) . Show that the \(\mathbf{S}_x\) cover \(\mathbf{X}\) , and that an arbitrary point of \(\mathbf{X}\) belongs only to a countable infinity of the \(\mathbf{S}_x\) {[}note that if \(y \in \mathbf{S}_x\) , we have \(d(x, y) \leqslant \frac{1}{2} r_y\) {]}.
\item
  Let R be the following equivalence relation between points \(x, y\) of X: there exists a sequence \((z_i)_{1 \leqslant i \leqslant n}\) of points of M such that \(x \in S_{z_i}\) , \(y \in S_{z_n}\) and \(S_{z_i}\) meets \(S_{z_{i+1}}\) for \(1 \leqslant i \leqslant n - 1\) . Show that the equivalence classes (mod R) are subspaces of X which are both open and closed in X and which have a countable base; in other words, X is the topological sum of metric spaces with countable bases (Chapter I, § 2, no. 4).
\end{enumerate}

¶ 19) In a metric space every relatively compact set is bounded. If X is a topological space, show that the following conditions on X are equivalent: α) there exists a metric on X, compatible with its topology, such that every bounded subset of X (with respect to this metric) is relatively compact; β) X is locally compact and has a countable base. {[}To show that α) → β), show that if every bounded set is relatively compact, then X is locally compact and σ-compact. To show that β) → α), note that X is metrizable by the corollary to Proposition 12; if d is a metric on X compatible with the topology of X, and if f is the function defined in § 1, Exercise 15 d), take the uniformity on X defined by the two pseudometrics d(x, y) and \textbar f(x) --- f(y)\textbar.{]}

¶ 20) a) Give an example of a closed equivalence relation R on a metrizable locally compact space X which has a countable base such that X/R is paracompact but such that there exists a point of X/R which has no countable fundamental system of neighbourhoods {[}cf.~Chapter I, § 10, Exercise 17 and Chapter IX, § 4, Exercise 24 a){]}.

\begin{enumerate}
\def\labelenumi{\alph{enumi})}
\setcounter{enumi}{1}
\tightlist
\item
  Let X be a metrizable space and let R be a closed equivalence relation on X. Show that if every point of X/R has a countable fundamental system of neighbourhoods, then each equivalence class (mod R) has a compact frontier in X. (Show that if the result were false there would exist a sequence in X, all of whose points were distinct, which had no cluster point, and whose image in X/R was a sequence which converged to a point distinct from all the points of the sequence.) For each \(z \in X/R\) , let \(C_{z}\) be the inverse image of z in X, and let \(F_{z}\) be the frontier of \(C_{z}\) if this frontier is not empty; if \(C_{z}\) is both open and closed, let \(F_{z}\) be any subset of \(C_{z}\) consisting of a single point. Let Y be the union of the sets \(F_{z}\) as z runs through X/R. Show that \(Y/R_{y}\) is homeomorphic to X/R.
\end{enumerate}

¶ 21) a) Let X be a metrizable space, let d be a bounded metric compatible with the topology of X, let σ be the metric corresponding to d on the set ℱ(X) of non-empty closed subsets of X (Exercise 6); and let R be an open and closed equivalence relation on X. Show that the restriction of σ to X/R is a metric compatible with the quotient topology. {[}Use Exercise 20 b) to show that every class mod R is open and compact. Then prove that if a sequence (zn) in X/R tends to a point a and if φ: X → X/R is the canonical mapping, we have σ(φ̄¹(a), φ̄¹(zₙ)) → 0; argue by contradiction, using the fact that R is open.{]}

\begin{enumerate}
\def\labelenumi{\alph{enumi})}
\setcounter{enumi}{1}
\tightlist
\item
  On the compact interval I = {[}0, 1{]} of R, let R be the equivalence relation whose classes are the points of the Cantor set K (Chapter VI, § 2, no. 5), other than the end-points of the intervals contiguous to K, and the closures of the intervals contiguous to K. Show that the quotient space I/R is homeomorphic to I {[}Chapter IV, § 8, Exercise 16 b){]} but that the metric \(\sigma\) is not compatible with the topology of I/R.
\end{enumerate}

\begin{enumerate}
\def\labelenumi{\arabic{enumi})}
\setcounter{enumi}{21}
\tightlist
\item
  Let X be a Hausdorff topological space with a countable base \((\mathbf{U}_{n})\) , and let R be a Hausdorff equivalence relation on X such that every point of X/R has a countable fundamental system of neighbourhoods. Show that the topology of X/R has a countable base. {[}Let \(\varphi: X \to X/R\) be the canonical mapping, and show that the interiors of finite unions of sets \(\varphi(\mathbf{U}_{n})\) form a base of the topology of X/R: if V is a neighbourhood of a point \(z \in X/R\) , and \((\mathbf{W}_{k})\) is a sequence of sets belonging to the base \((\mathbf{U}_{n})\) , contained in \(\overline{\varphi}^{1}(\mathbf{V})\) and covering \(\overline{\varphi}^{1}(z)\) , show that there is a finite number of indices k such that the union of the \(\varphi(\mathbf{W}_{k})\) is a neighbourhood of z. Argue by contradiction: if this statement were false we could construct a sequence \((y_{n})\) of distinct points of X/R, tending to z and such that \(y_{n}\) did not belong to the union of the \(\varphi(\mathbf{W}_{k})\) for \(k \leqslant n\) ; then show that the union of the \(\overline{\varphi}^{1}(y_{n})\) would be closed in X.{]}
\end{enumerate}

Show that the same conclusion is valid if R is assumed to be closed and every equivalence class mod R is compact (similar method).

¶ 23) A topological space is said to be submetrizable if its topology is finer than the topology of a metrizable space. A submetrizable space is Hausdorff but not necessarily regular (Chapter I, § 8, Exercise 20).

\begin{enumerate}
\def\labelenumi{\alph{enumi})}
\item
  A completely regular space X is submetrizable if and only if there exists a uniformity on X, compatible with the topology of X and defined by a family \(\Phi\) of pseudometrics which contains at least one metric d (cf.~§1, Exercise 3). Let \(\hat{X}\) be the completion of X with respect to this uniformity; the metric d extends to a pseudometric \(\overline{d}\) on \(\hat{X}\) ; show that for each \(x \in \hat{X}\) there is at most one point \(y \in X\) such that \(\overline{d}(x, y) = 0\) .
\item
  Show that if X is a completely regular submetrizable space, there exists a uniformity, compatible with the topology of X, and with respect to which X is complete {[}use a) and § 1, Exercise 16 c){]}. Hence show that if in addition X is pseudo-compact (§ 1, Exercise 21), then X is compact.
\item
  Show that a locally compact, paracompact, submetrizable space X is metrizable (reduce to the case where X is σ-compact by using Chapter I, § 9, no. 10, Theorem 5; then use the Corollary to Proposition 16) {[}cf.~§ 5, Exercise 15{]}.
\end{enumerate}

¶ 24) a) Let X be a complete metric space, and let A be a subset of X which is the intersection of a countable family of open sets. Show that there exists a metric on A with respect to which A is a complete metric space, and which defines on A the topology induced by that of X and also defines a uniformity finer than that induced by the uniformity of X {[}cf.~§ 1, Exercise 16 b) and c){]}.

\begin{enumerate}
\def\labelenumi{\alph{enumi})}
\setcounter{enumi}{1}
\tightlist
\item
  Conversely, let X be a metrizable space and let A be a subset of X such that there exists a metric d on A which is compatible with the topology induced by that of X and for which A is a complete metric space. Show that A is a countable intersection of open sets in X. {[}For each integer n \textgreater{} 0, consider the set \(G_{n}\) of points \(x \in \overline{A}\) which have an open neighbourhood U such that the diameter of \(A \cap U\) (with respect to d) is \(\leqslant r/n\) .{]}
\end{enumerate}

¶ 25) Let X be a metrizable space, let βX be its Stone-Čech compactification (§ 1, Exercise 7) and let d be a bounded metric compatible with the topology of X. For each x ∈ X, let f\_x(y) denote the real-valued function obtained by extending y → d(x, y) by continuity to βX.

\begin{enumerate}
\def\labelenumi{\alph{enumi})}
\item
  Show that \(f_{x}(z) + f_{y}(z) \geqslant d(x, y)\) whenever \(x\) and \(y\) are in \(\mathbf{X}\) and \(z\) is in \(\beta \mathbf{X}\) . For each \(y \in \beta \mathbf{X}\) other than \(x \in \beta \mathbf{X}\) , show that \(f_{x}(y) > 0\) . b) Show that if \(\mathbf{X}\) is complete with respect to the metric \(d\) , then \(\mathbf{X}\) is the intersection of a sequence of open sets in \(\beta \mathbf{X}\) . {[}Consider the set \(G_{n}\) of all \(y \in \beta \mathbf{X}\) such that \(f_{x}(y) < 1 / n\) for at least one point \(x \in \mathbf{X}\) , and show by using a) that \(\mathbf{X} = \bigcap_{n=1}^{\infty} G_{n}\) .{]}
\item
  Conversely, show that if X is the intersection of a sequence of open sets \(G_{n}\) in \(\beta X\) , there exists a metric \(d'\) on X which is compatible with the topology of X, and with respect to which X is complete. {[}Considering only the case in which \(X \neq \beta X\) , note that \(\beta X - X\) is the union of a family of compact sets \(F_{n} = \beta X - G_{n}\) ; for each n, let \(g_{n}(x) = \inf_{y \in F_{n}} f_{x}(y)\) , and show by using a) that \(g_{n}(x) > 0\) for all \(x \in X\) , and that \(g_{n}\) is continuous on X. Finish the proof as in Exercise 16 b) of § 1.{]}
\end{enumerate}

